\documentclass[12pt]{article}

\usepackage{sbc-template}
\usepackage[brazil]{babel}
\usepackage[utf8]{inputenc}
\usepackage[T1]{fontenc}
\usepackage{graphicx}
\usepackage{url}
\usepackage{booktabs}
\usepackage{amsmath}
\usepackage{float}
\usepackage{microtype}

\sloppy
\raggedbottom

% -------------------------------------------------------
\title{Análise de Sensibilidade Aplicada à Avaliação de Riscos\\
       em Portfólios de Investimentos}

\author{Anna J. A. Lucas\inst{1}, João V. M. Senra\inst{2}}


\address{
    Departamento de Ciência da Computação (DCC)\\
    Universidade Federal de Juiz de Fora (UFJF)\\
    Juiz de Fora -- MG -- Brasil\\
    \inst{1}\email{almeidaannajulia@gmail.com}, \inst{2}\email{joaovictor.senra@estudante.ufjf.br}
}

\begin{document}

\maketitle

% -------------------------------------------------------
\begin{abstract}
This paper presents a Decision Support System (DSS) for investment portfolio risk
analysis using Sensitivity Analysis. The system simulates seven market fluctuation
scenarios ranging from $-$20\% to $+$20\% and evaluates the financial impact on a
portfolio composed of eight Brazilian assets, including stocks, fixed income, and real
estate investment trusts (REITs). The application was developed in Python with an
interactive Streamlit dashboard. Results show that fixed-income assets act as a
stabilising buffer, limiting a 20\% market crash to a 14.8\% portfolio loss. The system
automatically generates allocation recommendations, supporting more informed investment
decisions.
\end{abstract}

\begin{resumo}
Este artigo apresenta um Sistema de Apoio à Decisão (SAD) para análise de risco de
portfólios de investimentos por meio de Análise de Sensibilidade. O sistema simula sete
cenários de oscilação de mercado de $-$20\% a $+$20\% e avalia o impacto financeiro
sobre uma carteira composta por oito ativos brasileiros, incluindo ações, renda fixa e
fundos imobiliários. A aplicação foi desenvolvida em Python com um painel interativo
Streamlit. Os resultados mostram que os ativos de renda fixa funcionam como um
amortecedor, limitando uma queda de 20\% no mercado a uma perda de 14,8\% na carteira.
O sistema gera recomendações de alocação automaticamente, apoiando decisões de
investimento mais conscientes.
\end{resumo}

% -------------------------------------------------------
\section{Introdução}

A gestão eficiente de um portfólio de investimentos exige ferramentas capazes de
antecipar os riscos associados às variações do mercado financeiro. Oscilações
macroeconômicas, crises setoriais e movimentos especulativos podem provocar perdas
expressivas em carteiras mal diversificadas, tornando a análise de risco um componente
essencial na tomada de decisão \cite{markowitz1952}.

Os Sistemas de Apoio à Decisão (SAD) são sistemas de informação projetados para
auxiliar tomadores de decisão em situações complexas e semiestruturadas \cite{power2002}.
No contexto financeiro, um SAD pode integrar modelos analíticos, bases de dados e
interfaces interativas para facilitar a compreensão do comportamento de uma carteira
diante de diferentes cenários de mercado.

A Análise de Sensibilidade é uma técnica de modelagem analítica que avalia como
variações nos parâmetros de entrada afetam os resultados de um modelo \cite{saltelli2008}.
Aplicada a portfólios de investimentos, ela possibilita simular cenários de alta e queda
nos preços dos ativos e mensurar o impacto dessas variações no valor total da carteira.
Essa abordagem corresponde à análise \textit{what-if} estruturada, amplamente utilizada
em ferramentas de suporte à decisão financeira \cite{damodaran2012}.

O \textbf{objetivo geral} deste trabalho é desenvolver um SAD interativo para análise
de risco de portfólios de investimentos por meio de Análise de Sensibilidade.
Os \textbf{objetivos específicos} são: \textit{(i)}~construir uma carteira simulada com
ativos financeiros brasileiros; \textit{(ii)}~implementar a análise de sensibilidade
para cenários de variação de $-$20\% a $+$20\%; \textit{(iii)}~calcular indicadores
de retorno e risco ponderados; e \textit{(iv)}~gerar recomendações automáticas de
alocação com base nos resultados obtidos.

% -------------------------------------------------------
\section{Metodologia}

\subsection{Base de Dados Simulada}

Foi construída uma base de dados sintética e realista, representando o portfólio de um
investidor pessoa física brasileiro. O arquivo \texttt{portfolio\_simulado.csv} contém
oito ativos com os campos: código, classe, quantidade, preço atual, retorno esperado
anual, volatilidade anualizada e perfil de risco. Os dados foram calibrados com base em
estimativas históricas de ativos equivalentes no mercado brasileiro em 2026, sem
dependência de conexão à internet ou APIs externas.

A Tabela~\ref{tab:portfolio} apresenta a composição da carteira. O valor total
investido é de R\$~67.863,00. Os ativos ITUB4 e VALE3 respondem por cerca de 30\% da
carteira, enquanto o Tesouro Selic representa 27,3\% , conferindo-lhe um importante
papel defensivo. A Figura~\ref{fig:alocacao} ilustra graficamente a distribuição.

\begin{table}[htbp]
\centering
\caption{Composição do Portfólio Simulado}
\label{tab:portfolio}
\begin{tabular}{llcrrr}
\toprule
\textbf{Ativo} & \textbf{Classe} & \textbf{Risco} &
\textbf{Ret. Esp.} & \textbf{Vol.} & \textbf{Participação} \\
\midrule
PETR4          & Ação       & Alto  & 18,0\% & 35,0\% & 11,3\% \\
VALE3          & Ação       & Alto  & 15,0\% & 30,0\% & 15,1\% \\
ITUB4          & Ação       & Médio & 12,0\% & 22,0\% & 15,4\% \\
BBDC4          & Ação       & Médio & 10,0\% & 20,0\% &  5,7\% \\
MGLU3          & Ação       & Alto  & 25,0\% & 55,0\% &  7,1\% \\
TESOURO SELIC  & Renda Fixa & Baixo & 10,7\% &  1,0\% & 27,3\% \\
HGLG11         & FII        & Médio &  9,0\% & 12,0\% & 12,0\% \\
XPML11         & FII        & Médio &  8,5\% & 14,0\% &  6,2\% \\
\bottomrule
\end{tabular}
\end{table}

\begin{figure}[htbp]
  \centering
  \includegraphics[width=0.55\linewidth]{fig_alocacao.png}
  \caption{Distribuição percentual do portfólio por ativo.}
  \label{fig:alocacao}
\end{figure}


\subsection{Cenários de Sensibilidade}

Foram definidos sete cenários predefinidos de variação de mercado, abrangendo desde
crises expressivas até rallies positivos, conforme a Tabela~\ref{tab:cenarios}.
O sistema também disponibiliza um \textit{slider} interativo que permite ao usuário
ajustar a variação manualmente de $-$30\% a $+$30\%.

\begin{table}[htbp]
\centering
\caption{Cenários de Análise de Sensibilidade}
\label{tab:cenarios}
\begin{tabular}{lc}
\toprule
\textbf{Cenário} & \textbf{Variação} \\
\midrule
Queda Forte      & $-20\%$ \\
Queda Moderada   & $-10\%$ \\
Queda Leve       & $-5\%$  \\
Neutro           &  $0\%$  \\
Alta Leve        & $+5\%$  \\
Alta Moderada    & $+10\%$ \\
Alta Forte       & $+20\%$ \\
\bottomrule
\end{tabular}
\end{table}

\subsection{Modelagem Matemática}

O preço simulado do ativo $i$ no cenário de variação $v$ é:
\begin{equation}
  P_i^{\text{sim}} = P_i^{\text{atual}} \times (1 + v \cdot \beta_i)
  \label{eq:preco}
\end{equation}
onde $\beta_i = 1{,}0$ para ações e FIIs e $\beta_i = 0{,}05$ para renda fixa,
refletindo a baixa correlação desses títulos com o mercado de capitais.

O impacto financeiro e o retorno do cenário são, respectivamente:
\begin{equation}
  \Delta_i = Q_i \!\cdot\! P_i^{\text{sim}} - V_i^{\text{ini}},
  \qquad
  R_{\text{cen}} = \frac{\sum \Delta_i}{\sum V_i^{\text{ini}}}
  \label{eq:impacto}
\end{equation}

O retorno esperado ponderado e o risco ponderado do portfólio são:
\begin{equation}
  \bar{R} = \sum_{i} w_i r_i = 13{,}0\%~\text{a.a.},
  \qquad
  \sigma_p = \sum_{i} w_i \sigma_i = 19{,}5\%~\text{a.a.}
  \label{eq:ponderados}
\end{equation}

\subsection{Implementação}

O sistema foi desenvolvido em Python~3 com as bibliotecas Pandas \cite{mckinney2010}
(manipulação de dados), NumPy (operações vetoriais), Plotly (gráficos interativos) e
Streamlit (interface web). Toda a lógica é executada localmente, sem dependência de
credenciais ou conexão à internet, sendo compatível com Windows, macOS e Linux.

% -------------------------------------------------------
\section{Resultados e Discussões}

\subsection{Indicadores Globais}

O portfólio simulado possui valor total de R\$~67.863,00, retorno esperado ponderado
de 13,0\% a.a.\ e volatilidade ponderada de 19,5\% a.a. O Tesouro Selic (27,3\% da
carteira) atua como âncora defensiva, enquanto o conjunto de ações representa 54,5\%
do portfólio.

\subsection{Comportamento nos Cenários Simulados}

A Tabela~\ref{tab:resultados} apresenta o impacto financeiro em cada cenário, e a
Figura~\ref{fig:sensibilidade} ilustra a curva de sensibilidade do portfólio para
variações de $-$30\% a $+$30\%.

\begin{table}[htbp]
\centering
\caption{Impacto Financeiro nos Cenários de Mercado}
\label{tab:resultados}
\begin{tabular}{lrr}
\toprule
\textbf{Cenário} & \textbf{Ganho/Perda (R\$)} & \textbf{Retorno} \\
\midrule
Queda Forte ($-$20\%)    & $-$R\$~10.058 & $-$14,8\% \\
Queda Moderada ($-$10\%) & $-$R\$~5.029  & $-$7,4\%  \\
Queda Leve ($-$5\%)      & $-$R\$~2.514  & $-$3,7\%  \\
Neutro (0\%)              &  R\$~0        &  0,0\%    \\
Alta Leve ($+$5\%)        & $+$R\$~2.514  & $+$3,7\%  \\
Alta Moderada ($+$10\%)   & $+$R\$~5.029  & $+$7,4\%  \\
Alta Forte ($+$20\%)      & $+$R\$~10.058 & $+$14,8\% \\
\bottomrule
\end{tabular}
\end{table}

\begin{figure}[htbp]
  \centering
  \includegraphics[width=0.95\linewidth]{fig_sensibilidade.png}
  \caption{Curva de sensibilidade: valor do portfólio em função da variação de mercado.
           A linha tracejada indica o valor inicial (R\$~67.863). Pontos marcados
           correspondem aos cenários predefinidos.}
  \label{fig:sensibilidade}
\end{figure}

Observa-se uma relação aproximadamente linear entre a variação de mercado e o impacto
na carteira, efeito esperado da modelagem proporcional adotada. O papel amortecedor do
Tesouro Selic é evidente: uma queda de 20\% no mercado de renda variável resulta em
apenas 14,8\% de perda no portfólio total, em vez dos 20\% que ocorreriam em uma
carteira integralmente composta por ações.

\subsection{Impacto por Ativo}

A Figura~\ref{fig:impacto} detalha o impacto financeiro individual no cenário de Queda
Moderada ($-$10\%). ITUB4 e VALE3 registram as maiores perdas absolutas em razão de sua
elevada participação na carteira. O MGLU3, apesar da maior volatilidade esperada (55\%),
apresenta impacto absoluto menor devido ao preço unitário reduzido. O Tesouro Selic
registra perda mínima ($\approx -$R\$~93), confirmando seu papel defensivo.

\begin{figure}[htbp]
  \centering
  \includegraphics[width=0.95\linewidth]{fig_impacto.png}
  \caption{Impacto financeiro por ativo no cenário de Queda Moderada ($-$10\%).}
  \label{fig:impacto}
\end{figure}

\subsection{Recomendações de Alocação}

Com base nos resultados, o SAD gerou as seguintes recomendações automáticas:

\begin{itemize}
  \item \textbf{Manter a âncora de renda fixa:} o Tesouro Selic limita
    significativamente as perdas em cenários adversos. Ampliar sua participação para
    acima de 30\% aumentaria a resiliência do portfólio.
  \item \textbf{Revisar a exposição ao MGLU3:} com volatilidade de 55\% a.a., esse
    ativo eleva o risco ponderado do portfólio sem contribuição proporcional no valor
    investido.
  \item \textbf{Diversificação setorial adequada:} a presença de ativos de energia,
    mineração, setor financeiro, varejo e imobiliário reduz o risco idiossincrático.
  \item \textbf{Perfil de risco coerente:} 54,5\% da carteira em renda variável é
    compatível com um perfil moderado, mas exige monitoramento em cenários de queda
    acima de 10\%.
\end{itemize}

A limitação central do modelo é a não consideração da correlação entre ativos. O risco
real de um portfólio diversificado é geralmente inferior à média ponderada das
volatilidades individuais \cite{bodie2014}, o que indica que as estimativas aqui
apresentadas tendem a ser conservadoras.

% -------------------------------------------------------
\section{Conclusão}

Este trabalho desenvolveu um SAD interativo para análise de risco em portfólios de
investimentos por meio de Análise de Sensibilidade. O sistema simula sete cenários de
oscilação de mercado, calcula o impacto individual por ativo e no portfólio total, e
gera recomendações automáticas de alocação. Os resultados evidenciam que a presença de
ativos de renda fixa atua como fator de estabilização, reduzindo a sensibilidade da
carteira a movimentos adversos do mercado.

O SAD desenvolvido cumpre o objetivo de tornar visível e quantificável o risco
associado a diferentes composições de carteira, apoiando decisões mais informadas sobre
alocação de ativos. A interface interativa em Streamlit facilita a exploração dos
cenários por usuários sem necessidade de conhecimento técnico em programação.

Como limitações, destaca-se a ausência da correlação entre ativos no cálculo do risco
(matriz de covariâncias de Markowitz), o uso de betas simplificados e a base de dados
sintética. Trabalhos futuros poderão incorporar a Teoria Moderna de Portfólio
\cite{markowitz1952}, simulação de Monte Carlo e o cálculo do \textit{Value at Risk}
(VaR) para uma análise mais precisa do risco de portfólio.

% -------------------------------------------------------
\begin{thebibliography}{6}

\bibitem{markowitz1952}
MARKOWITZ, H.
Portfolio selection.
\textit{The Journal of Finance}, v.~7, n.~1, p.~77--91, 1952.

\bibitem{power2002}
POWER, D.~J.
\textit{Decision Support Systems: Concepts and Resources for Managers}.
Westport: Quorum Books, 2002.

\bibitem{saltelli2008}
SALTELLI, A. et al.
\textit{Global Sensitivity Analysis: The Primer}.
Chichester: John Wiley \& Sons, 2008.

\bibitem{damodaran2012}
DAMODARAN, A.
\textit{Investment Valuation: Tools and Techniques for Determining the Value of
any Asset}. 3.~ed. Hoboken: Wiley Finance, 2012.

\bibitem{mckinney2010}
McKINNEY, W.
Data structures for statistical computing in Python.
In: \textit{Proceedings of the 9th Python in Science Conference}, 2010.
p.~56--61.

\bibitem{bodie2014}
BODIE, Z.; KANE, A.; MARCUS, A.~J.
\textit{Investments}. 10.~ed. New York: McGraw-Hill Education, 2014.

\end{thebibliography}

\end{document}
